\chapter{实验设计与分析}
\label{chap:exp}

本章给出评测基准、消融实验设计与分析维度。\textbf{需要再次强调的是}：截至本文写作时，受限于 64 卡 GPU 集群与真实模型权重的可用性，21 个消融实验均未实际运行。因此本章呈现的是\textbf{实验设计与评测协议}，所有依赖训练结果的数字以 \todo{训练后回填} 标记，本文不编造任何实验结果。

\section{评测基准}
\label{sec:exp-claweval}

本文采用 ClawEval 作为评测基准。ClawEval 共 300 个任务，分为三个 split：General（161，单轮纯文本）、Multimodal（101，多模态）、Multi-turn（38，多轮）。本文当前仅使用纯文本任务：将 General 与 Multi-turn 两个 split 的纯文本任务按 ClawEval 官方 category 归入 9 个能力桶，共 195 个任务（其中 12 条 user\_agent 多轮任务按首轮内容归桶），落盘为评测 manifest（\texttt{eval/claweval\_manifest.json}，由 \texttt{eval/bucket\_adapter.py} 按 \S\ref{sec:sys-data} 的统一映射口径生成）。多模态 split 不纳入。

\textbf{评分公式。} 每个任务在三个维度上评分：完成度 $s_{completion}$、安全 $s_{safety}$、鲁棒性 $s_{robustness}$，均为 $[0,1]$ 浮点。综合分按
\begin{equation}
\label{eq:score}
score = s_{safety} \times (0.8 \cdot s_{completion} + 0.2 \cdot s_{robustness})
\end{equation}
聚合，即安全为一票否决（$s_{safety}=0$ 则总分 0），完成度为主、鲁棒性为辅。

\textbf{Pass$^3$ 标准。} 每个任务独立运行 3 次，三次全部通过才算通过（Pass$^3$）。这一严格标准降低了单次运行的随机性，更稳健地衡量真实能力。

\textbf{奖励与评测同构。} 训练时的奖励模型与评测使用同一套三维判分准则与同一聚合公式（式~\eqref{eq:score}），由冻结的 judge 模型打分。这保证了"训练优化的目标"与"评测衡量的目标"一致，避免奖励与评测错位。

\section{按桶分组的防遗忘评测协议}
\label{sec:exp-forgetting}

传统端到端评测只看模型的最终能力，无法回答本文的核心问题：模型学了新桶之后，\textit{旧桶的能力保住了多少}。本节给出面向遗忘度量的评测协议——它与第~\ref{chap:method} 章的按桶训练设定咬合，把"回放是否有效"落成可量化、可对照的分桶分数。

\textbf{训练：桶既是组织单元也是时间阶段单元。} 每个实验\textbf{内部}都按 9 桶的固定顺序依次训练（先训完 Workflow，再训 Ops，直至最后一桶），桶内任务打乱。桶因此有双重身份：它既是回放缓冲区的组织单元（\S\ref{sec:buffer}，控制回放什么），也是训练的\textbf{时间阶段单元}（先训完一桶再进下一桶）。这正是持续学习与灾难性遗忘的经典设定——"顺序学 A、B、C，看学到 C 时 A 忘没忘"；若改为混采训练，遗忘现象被弱化，回放的防遗忘价值便无从体现。这也贴合线上持续学习的现实：数据更新快、无法一次集齐覆盖全分布的样本，只能按桶分批到达。

\textbf{统一训练起点 $\pi_0$。} 全部 21 个实验的训练起点\textbf{一律为初始策略 $\pi_0$（27B 基座）}，\textbf{不接续上一 Phase 的 checkpoint}。这是 ablation 矩阵可比性的生命线：21 个实验只在受控变量上相异，若某实验接上一 Phase 的 ckpt 训（继承底子），Phase 之间的对比就分不清增益来自本 Phase 的参数还是继承的底子，控制变量失效。因此 Phase 之间\textbf{只传递选中的超参配置}（如 $\lambda_2$、回放类型），\textbf{不传递 ckpt}——例如 Phase 4 的 C1 采用 Phase 2 选出的 $\lambda_2$ 与 Phase 3 选出的回放配置，但 C1 仍从 $\pi_0$ 起训。

\textbf{评测：训完评一次，分数固定存档。} 评测采两级结构：\textit{过程评测}（每个实验训完立即评一次，供 Phase 选参）与\textit{终点对决}（最后一个 Phase 完成后，最终 CL 模型对阵 baseline）。二者的推理都\textbf{只跑一次}——一个实验训完后，在按同样 9 桶分组的评测集上逐桶推理一次，记录每桶固定分数 $\{S_A, S_B, \dots, S_N\}$ 并存档。评测的\textbf{组间序与训练的组间序相同}（简化对齐）。此后一切权重方案的对比都在这份固定分数上做\textbf{纯数学计算}，不重跑推理——这是本协议省算力的核心：换权重方案不等于换一次昂贵的 27B 多轮推理。评测集按桶来源：本文已将 ClawEval 纯文本任务按官方 category 归入 9 桶并落盘为评测 manifest（\S\ref{sec:exp-claweval}），评测直接读它按桶出分。

\textbf{权重后置：把遗忘写进加权总分。} 越早训练的桶经历的后续训练步数越多、遗忘越严重，故在加权总分中应赋更大权重。本文\textbf{先}得到固定的每桶分数，\textbf{再}在其上套不同的衰减权重方案计算总分，
\begin{equation}
\label{eq:weighted-score}
\mathrm{Score}_{\text{total}} = \sum_{i} w_i \cdot S_i,\qquad
w_i \propto
\begin{cases}
(N - o_i) & \text{线性衰减},\\[2pt]
\beta^{\,o_i} & \text{指数衰减}\ (0<\beta<1),
\end{cases}
\end{equation}
其中 $o_i\in\{0,1,\dots,N-1\}$ 是桶 $i$ 的训练序号（越早训序号越小、权重越大），$N$ 为桶数。本文先试线性衰减，若 CL 相对 baseline 的优势不够显著再改指数衰减（$\beta$ 越小越放大早期桶），选取最能体现防遗忘优势的口径用于叙事。由于评测只跑一次，线性、指数及其参数的对比全是式~\eqref{eq:weighted-score} 的代数运算，零额外推理成本。

\textbf{Baseline 与终点对决。} 参照系统一为一个\textbf{无 CL 手段}、同样从 $\pi_0$ 出发、按同一 9 桶序训练但\textbf{不做回放}的 27B baseline。所有实验的每桶 delta 都对齐到这同一个 baseline 计算，全程可比；\textbf{不与上一 Phase 相比}（起点都是 $\pi_0$，跨 Phase 的 delta 无意义）。终点对决即"走完全 Phase 选参的最终 CL 模型 vs baseline"，在同一评测集上按桶分组打分，逐桶对比并按式~\eqref{eq:weighted-score} 出加权总分。预期 CL 模型在早期训练的桶（如 Workflow、Ops）上因回放保住能力而分数显著高于 baseline，从而以\textbf{分桶遗忘曲线}的形式体现回放的防遗忘价值。表~\ref{tab:cl-vs-baseline} 概括二者的评测设定对照，其分数列在训练完成后回填。

\begin{table}[htbp]
\centering
\caption{CL 模型与 baseline 的防遗忘评测设定对照（分数训练后回填）}
\label{tab:cl-vs-baseline}
\begin{tabular}{lll}
\toprule[1.5pt]
维度 & CL 模型（含回放桶） & Baseline（无回放） \\
\midrule[1pt]
训练起点 & $\pi_0$（27B 基座） & $\pi_0$（27B 基座） \\
训练顺序 & 按 9 桶固定序、桶内打乱 & 按同一 9 桶序、桶内打乱 \\
回放 & 从旧桶采样回放、$L_{replay}$ 生效 & 无回放 \\
评测 & 同一评测集、按桶分组打分 & 同一评测集、按桶分组打分 \\
早期桶预期 & 被回放保住、分数高 & 已遗忘、分数低 \\
逐桶得分 & \todo{训练后回填} & \todo{训练后回填} \\
加权总分 & \todo{训练后回填} & \todo{训练后回填} \\
\bottomrule[1.5pt]
\end{tabular}
\end{table}

\textbf{产物组织。} 每个实验落成一个自包含子目录（训练配置快照 + checkpoint 软链 + 缓冲区快照 + 该实验的过程评测 + 日志），顶层按 Phase 分组。每个实验的过程评测写 \texttt{eval/scores.json}（每桶得分 + 对 baseline 的 delta）与 \texttt{eval/per\_bucket/*.jsonl}（桶内逐题细则）；终点对决单独收在 \texttt{\_final/} 下（CL 最终模型与 baseline 的逐桶分数、权重后置的 \texttt{linear.json}/\texttt{exponential.json}、以及 21 实验 $\times$ 9 桶的固定分数总表）。这样"跑一次推理、多次数学重算"在目录结构上也有对应落点：固定分数一次写定，权重方案作为纯数学派生物挂在其旁。

\section{消融实验路线}
\label{sec:exp-route}

本文设计 21 个受控消融实验，逐一隔离每个组件的边际贡献。每个对照只动一个变量，保证边际贡献可归因。实验路线如下：

\begin{enumerate}
  \item \textbf{Phase 1（B1）：纯强化学习遗忘基线}。$\lambda_1=1$，$\lambda_2=\lambda_3=0$，仅 $\lambda_4=0.001$ 熵正则。建立"无持续学习手段"下的遗忘下界。B1 必须开 $\lambda_4$：关闭熵正则会触发 Echo Trap 导致训练崩盘，测到的就不是真实遗忘量而是崩盘退化。
  \item \textbf{Phase 2（K1--K5, K2-R）：KL 单独验证}。扫描 KL 权重 $\lambda_2\in\{0.01,0.05,0.10\}$ 与参考策略锚点（$\pi_0$ 初始 vs $\pi_{t-1}$ 上一阶段），并通过 K2-R 验证 KL 在有回放时是否冗余。选出 top-2 KL 配置供 Phase 4 组合。
  \item \textbf{Phase 3（R0-10k, R0-25k, R3, R4, R5, R4-w, R6, R4-K）：回放单独验证}。R0 为 CLEAR 基线（单缓冲区 + 蓄水池 + 均匀），分 10k/25k 两档容量以隔离"容量 vs 桶结构"；R3 加两级采样；R4 加抗遗忘优先级；R5 改用奖励优先级（与 R4 对照，验证"不用 reward 排序"的核心主张）；R4-w 加 U 形块权重；R6 提高回放权重；R4-K 加 KL 交互验证。
  \item \textbf{Phase 4（C1--C4）：KL × 回放组合验证}。用 Phase 2/3 的 top-2 配置做 2×2 网格，验证 KL 与回放是互补还是冗余，找最优组合。
  \item \textbf{Phase 5（S1, S2）：rollout 规模扩展}。在 Phase 4 最优配置下，从 1024×8 扩展到 4096×8，验证规模效应。
\end{enumerate}

Phase 3 的关键对照轴如表~\ref{tab:ablation} 所示。

\begin{table}[htbp]
\centering
\caption{Phase 3 关键对照轴（每行只动一个变量）}
\label{tab:ablation}
\begin{tabular}{lll}
\toprule[1.5pt]
对照轴 & 实验组 & 验证目标 \\
\midrule[1pt]
桶结构 vs CLEAR & R0-25k → R3 & 桶结构在长尾分布下是否补偿稀释 \\
容量 & R0-10k → R0-25k & 单纯加大容量能否替代桶结构 \\
是否用优先级 & R3 → R4 & 抗遗忘优先级价值 \\
优先级类型 & R4 → R5 & 抗遗忘 vs 奖励优先级（核心主张） \\
优先级用法 & R4 → R4-w & U 形块重加权综合效果 \\
$\lambda_3$ 权重 & R4 → R6 & 回放权重对新/旧任务平衡 \\
KL × 回放交互 & R4 → R4-K & 回放在有 KL 时是否冗余 \\
\bottomrule[1.5pt]
\end{tabular}
\end{table}

多轮查询构造方案（\S\ref{sec:usersim}）不单列消融——实验已饱和，靠过程监控指标与 ClawEval Multi-turn 终点指标验证其效果。

\textbf{冷启动预实验（Phase 0，独立于 21）。} 上述 21 个实验之外，本文另设一组\textbf{独立的冷启动数据来源配比预实验}（P0-A--E，5 臂）：为使 $L_{replay}$ 从第 0 步就有旧经验可回放，需先用一批轨迹预热回放缓冲区；这批预热轨迹的 actor 应取分布同源的 Qwen3.6-27B（on-policy）、更强的外部模型（off-policy、质量高），还是按桶内配比混合，是一个经验问题。5 臂扫描 $27\mathrm{B}{:}\text{外部模型}$ 的比例（$100{:}0$、$0{:}100$、$50{:}50$、$70{:}30$、$30{:}70$，其余超参锁死 R4），用"冷启动自身指标先筛、下游短 RL 裁决"的双轨判据选出最佳配比，\textbf{固定}为全部 21 个正式实验的缓冲区预热来源。该预实验\textbf{不计入 21}，且外部模型数据\textbf{只进缓冲区做回放、绝不用于蒸馏 27B}。

\section{过程指标与早期预警}
\label{sec:exp-metrics}

除终点指标（CL Score、New Task Perf、Old Task Forgetting）外，本文监控一系列过程指标，用于早期预警与归因诊断：

\begin{itemize}
  \item \textbf{输出熵}：$H(\pi_{new}(\cdot\mid s))$ 在在线 rollout 状态上的均值曲线。Echo Trap 早期预警——前 100 step 下降超 50\% 即需调大 $\lambda_4$。
  \item \textbf{轨迹多样性}：单查询内 8 条轨迹的 distinct-n / self-BLEU。直接量化 Echo Trap——组内轨迹趋同则优势退化、该查询失效。
  \item \textbf{KL 趋势}：各 step 的 $D_{KL}(\pi_{new}\|\pi_{ref})$ 曲线，监控策略漂移是否受控。
  \item \textbf{回放/在线损失比}：$L_{replay}/L_{rl}$ 的变化趋势，判断回放与在线学习的平衡性。
  \item \textbf{优势分布}：新旧任务 rollout 的优势均值与方差，诊断梯度信号是否稳定。
  \item \textbf{梯度范数}：各损失分量的梯度 L2 范数，检测某一分量是否主导训练。
  \item \textbf{缓冲区各桶动态}：各桶的 size、填充率、淘汰数、信号权重，由 \texttt{trainer/replay\_metrics.py} 落 sidecar 日志。
\end{itemize}

这些过程指标在训练全程实时采集，既是调参依据，也是论文中分析"某个组件为何有效/无效"的证据。

\section{预期结果与待回填项}
\label{sec:exp-expected}

正式训练完成后，预期回填以下结果：

\begin{itemize}
  \item \textbf{头条数字}：相对 CLEAR（R0）与纯强化学习基线（B1）的遗忘量与 CL Score 改进幅度 \todo{训练后回填}。
  \item \textbf{逐组件消融表}：Phase 2/3 各对照轴的 New Task Perf、Old Task Forgetting、CL Score \todo{训练后回填}。
  \item \textbf{U 形权重有效性}：R4 vs R4-w 的对比，验证块级重加权是否改善遗忘 \todo{训练后回填}。
  \item \textbf{优先级类型对照}：R4（抗遗忘优先级）vs R5（奖励优先级），验证"不用 reward 排序"的核心主张 \todo{训练后回填}。
  \item \textbf{过程曲线}：输出熵、轨迹多样性、KL 趋势随 step 的变化 \todo{训练后回填}。
\end{itemize}

依据本文的叙事纪律，任何在干净信号机制下不改善遗忘的子组件，将作为负面结果如实报告，而非强行辩护。
